\documentclass[../../main.tex]{subfiles}
\begin{document}

\begin{center}
  {\fontsize{48pt}{58pt}\selectfont \textbf{1999}\normalsize\textbf{年度}}
\end{center}

\begin{table}[htb]
  \centering
  \caption{出題テーマと難易度}
  \begin{tabular}{|c|p{9cm}|c|}
    \hline
    \textbf{問題番号} & \textbf{テーマ} & \textbf{難易度} \\
    \hline
    第1問 & 絶対値を含む曲線と直線が囲む面積比 & \\
    \hline
    第2問 & 三角形の角と正弦の積の最大値 & \\
    \hline
    第3問 & 非線形漸化式と数列の比 & \\
    \hline
    第4問 & 合同な面からなる四面体の構成 & \\
    \hline
    第5問 & 二次方程式の解と整数性 & \\
    \hline
    第6問 & 積分の平均値と回転体の分割 & \\
    \hline
  \end{tabular}
\end{table}

\begin{table}[htb]
  \centering
  \caption{解答}
  \renewcommand{\arraystretch}{1.4}
  \begin{tabular}{|c|c|p{9cm}|}
    \hline
    \textbf{問題番号} & \textbf{小問} & \textbf{答え} \\
    \hline
    第1問 & -- & 傾き$\dfrac{1}{2}$ \\
    \hline
    第2問 & -- & 最大値$\dfrac{3\sqrt{3}}{8}$ \\
    \hline
    第3問 & (1) & $(a_2,b_2)=(2\alpha,-2\alpha)$，$(a_3,b_3)=(-2\alpha^2,-2\alpha^2)$，$(a_4,b_4)=(8\alpha^5,-8\alpha^5)$ \\
    \cline{2-3}
    & (2) & $\dfrac{a_{2n+1}}{a_{2n}}=-\alpha$ \\
    \hline
    第4問 & -- & 証明問題 \\
    \hline
    第5問 & -- & 証明問題 \\
    \hline
    第6問 & (1) & 証明問題 \\
    \cline{2-3}
    & (2) & $1-\dfrac{8}{\pi^2}$ \\
    \hline
  \end{tabular}
\end{table}

\paragraph{試験形態}
6問．試験時間は要確認．

\end{document}
